% ============================================================
% Pipeline básico de um projeto de IA — Apresentação didática
% Autor: Alexandre Barros (alexand7e.dev.br)
% Repo : https://github.com/alexand7e/Pipeline-IA-do-Zero
% ============================================================
\documentclass[aspectratio=169]{beamer}

\usetheme{Madrid}
\usecolortheme{default}

\usepackage[T1]{fontenc}
\usepackage[utf8]{inputenc}
\usepackage[brazil]{babel}
\usepackage{amsmath, amssymb}
\usepackage{booktabs}
\usepackage{array}

% ── identidade visual ─────────────────────────────────────
\setbeamertemplate{navigation symbols}{}
\setbeamertemplate{footline}{
  \leavevmode\hbox{%
    \begin{beamercolorbox}[wd=.45\paperwidth,ht=2.25ex,dp=1ex,left,leftskip=2ex]{author in head/foot}
      \usebeamerfont{author in head/foot}Alexandre Barros \textbullet{} alexand7e.dev.br
    \end{beamercolorbox}%
    \begin{beamercolorbox}[wd=.45\paperwidth,ht=2.25ex,dp=1ex,center]{title in head/foot}
      \usebeamerfont{title in head/foot}Pipeline IA do Zero
    \end{beamercolorbox}%
    \begin{beamercolorbox}[wd=.1\paperwidth,ht=2.25ex,dp=1ex,right,rightskip=2ex]{date in head/foot}
      \usebeamerfont{date in head/foot}\insertframenumber/\inserttotalframenumber
    \end{beamercolorbox}}%
  \vskip0pt%
}

\title[Pipeline IA do Zero]{Pipeline básico de um projeto de IA}
\subtitle{Do problema real à decisão em produção}
\author{Alexandre Barros}
\institute{%
  \href{https://alexand7e.dev.br}{alexand7e.dev.br}\\[4pt]
  \footnotesize\texttt{github.com/alexand7e/Pipeline-IA-do-Zero}
}
\date{\today}

% ============================================================
\begin{document}
% ============================================================

% ── CAPA ────────────────────────────────────────────────────
\begin{frame}
  \titlepage
\end{frame}

% ── OBJETIVO ────────────────────────────────────────────────
\begin{frame}{O que vamos ver hoje}
\begin{itemize}
  \item As \textbf{7 etapas} que todo projeto de IA percorre — do dado à predição.
  \item Como cada escolha afeta a capacidade do modelo de \textbf{generalizar}.
  \item Os problemas clássicos: \textbf{overfitting, underfitting, viés e variância}.
  \item Como \textbf{avaliar de forma honesta} — além da acurácia.
\end{itemize}
\vfill
\begin{block}{Proposta}
  Entender os \emph{porquês} antes de usar frameworks. \
  Um modelo que você entende é um modelo que você consegue corrigir.
\end{block}
\end{frame}

% ── PROBLEMA REAL ───────────────────────────────────────────
\begin{frame}{O problema: classificação binária}
\begin{columns}
\column{0.55\textwidth}
\textbf{Exemplos reais:}
\begin{itemize}
  \item Fraude \textbf{vs} não fraude
  \item Churn \textbf{vs} fidelidade
  \item Diagnóstico \textbf{vs} saudável
  \item Aprovado \textbf{vs} reprovado
\end{itemize}
\vspace{8pt}
Em todos os casos:
\begin{itemize}
  \item Temos \textbf{features} $X$ (sinais observados)
  \item Queremos prever um \textbf{rótulo} $y \in \{0, 1\}$
  \item O modelo devolve uma \textbf{probabilidade} $\hat{p} \in (0, 1)$
\end{itemize}
\column{0.4\textwidth}
\begin{block}{Decisão limiar}
  \[
    \hat{y} =
    \begin{cases}
      1 & \text{se } \hat{p} \geq 0{,}5 \\
      0 & \text{caso contrário}
    \end{cases}
  \]
  O limiar 0,5 é ajustável — depende do custo de cada tipo de erro.
\end{block}
\end{columns}
\end{frame}

% ── MAPA DOS 7 PASSOS ───────────────────────────────────────
\begin{frame}{O pipeline em 7 passos}
\begin{center}
\begin{tabular}{cl}
  \toprule
  \textbf{Passo} & \textbf{O que acontece} \\
  \midrule
  1 & \textbf{Definir} o problema e escolher a(s) métrica(s) de sucesso \\
  2 & \textbf{Obter} dados representativos e reproduzíveis \\
  3 & \textbf{Preparar} os dados (transformar, normalizar) \\
  4 & \textbf{Dividir} em treino, validação e teste \\
  5 & \textbf{Treinar} o modelo — aprender pesos a partir dos dados \\
  6 & \textbf{Avaliar} — medir generalização e analisar erros \\
  7 & \textbf{Empacotar} — levar o modelo para produção \\
  \bottomrule
\end{tabular}
\end{center}
\vfill
\begin{alertblock}{Importante}
  Cada passo influencia todos os seguintes.
  Pular ou errar o passo 1 compromete os 6 restantes.
\end{alertblock}
\end{frame}

% ── PASSO 1: PROBLEMA E MÉTRICA ─────────────────────────────
\begin{frame}{Passo 1 — Definir o problema e a métrica}
\begin{columns}
\column{0.5\textwidth}
\textbf{Perguntas obrigatórias:}
\begin{itemize}
  \item O que queremos prever?
  \item Como saberemos se acertamos?
  \item Qual erro é mais caro — falso positivo ou falso negativo?
\end{itemize}
\vspace{6pt}
\textbf{Métricas comuns:}
\begin{itemize}
  \item Acurácia, Precisão, Recall, F1
  \item AUC-ROC (robusta ao desbalanceamento)
\end{itemize}
\column{0.45\textwidth}
\begin{alertblock}{Acurácia pode enganar}
  Dataset com 95\% de classe 0:\\
  um modelo que sempre prediz 0\\
  tem \textbf{95\% de acurácia} e \textbf{recall = 0}.\\[4pt]
  \textit{Métricas definem o que "bom" significa — escolha antes de treinar.}
\end{alertblock}
\end{columns}
\end{frame}

% ── MÉTRICAS ────────────────────────────────────────────────
\begin{frame}{Métricas de classificação binária}
\begin{columns}
\column{0.5\textwidth}
\[
  \text{Precisão} = \frac{TP}{TP + FP}
  \quad \text{(confiança no positivo)}
\]
\[
  \text{Recall} = \frac{TP}{TP + FN}
  \quad \text{(cobertura do positivo)}
\]
\[
  F_1 = \frac{2 \cdot P \cdot R}{P + R}
  \quad \text{(balanço)}
\]
\column{0.45\textwidth}
\begin{center}
\textbf{Matriz de confusão}\\[6pt]
\begin{tabular}{r|cc}
         & Pred 0 & Pred 1 \\
\hline
Real 0   & \textcolor{blue}{TN} & \textcolor{red}{FP} \\
Real 1   & \textcolor{red}{FN} & \textcolor{blue}{TP} \\
\end{tabular}
\vspace{8pt}

\footnotesize
\textcolor{red}{FP} = alarme falso \\
\textcolor{red}{FN} = caso perdido \\[4pt]
\textit{O custo de cada erro depende do domínio.}
\end{center}
\end{columns}
\end{frame}

% ── PASSO 2: DADOS ──────────────────────────────────────────
\begin{frame}{Passo 2 — Obter dados}
\begin{columns}
\column{0.55\textwidth}
\begin{block}{Princípio fundamental}
  \textit{``Garbage in, garbage out''} \\
  Um modelo é tão bom quanto os dados que o treinaram.
\end{block}
\vspace{8pt}
\textbf{O que garantir:}
\begin{itemize}
  \item \textbf{Representatividade} — os dados refletem o mundo real?
  \item \textbf{Reprodutibilidade} — semente aleatória fixa (\textit{seed})
  \item \textbf{Qualidade} — valores ausentes, outliers, inconsistências
  \item \textbf{Volume} — suficiente para generalizar e avaliar
\end{itemize}
\column{0.4\textwidth}
\begin{exampleblock}{Neste projeto}
  Dataset sintético \textit{``duas luas''}: \\
  problema não-linear com 2 features, \\
  visualmente interpretável.\\[4pt]
  Serve para focar no \emph{pipeline}, \\
  não nas particularidades dos dados.
\end{exampleblock}
\end{columns}
\end{frame}

% ── PASSO 3: PRÉ-PROCESSAMENTO ──────────────────────────────
\begin{frame}{Passo 3 — Preparar os dados}
\begin{columns}
\column{0.5\textwidth}
\textbf{Padronização Z-score:}
\[
  X' = \frac{X - \mu}{\sigma}
\]
\begin{itemize}
  \item Coloca todas as features na mesma escala
  \item Melhora convergência do gradiente
\end{itemize}
\vspace{8pt}
\textbf{Regra de ouro — evitar vazamento:}
\begin{itemize}
  \item $\mu$ e $\sigma$ calculados \textbf{só no treino}
  \item Aplicados igualmente em validação, teste e produção
\end{itemize}
\column{0.45\textwidth}
\begin{alertblock}{Vazamento de dados (\textit{data leakage})}
  Calcular a normalização com teste/validação é usar ``informação do futuro'' no treino.\\[4pt]
  Resulta em estimativas de performance otimistas — e surpresas em produção.
\end{alertblock}
\end{columns}
\end{frame}

% ── PASSO 4: SPLIT ──────────────────────────────────────────
\begin{frame}{Passo 4 — Dividir os dados}
\begin{center}
\begin{tabular}{lcl}
  \toprule
  \textbf{Conjunto} & \textbf{Proporção} & \textbf{Papel} \\
  \midrule
  Treino     & 70\% & Ajusta os pesos do modelo \\
  Validação  & 15\% & Guia decisões durante o desenvolvimento \\
  Teste      & 15\% & Avalia a performance \textbf{final} (usar uma só vez) \\
  \bottomrule
\end{tabular}
\end{center}
\vfill
\begin{columns}
\column{0.48\textwidth}
\begin{block}{Por que validação separada?}
  Cada vez que você toma uma decisão olhando para o teste (escolher arquitetura, hiperparâmetros), você indiretamente treina nele.
  A validação existe para proteger a integridade do teste.
\end{block}
\column{0.48\textwidth}
\begin{block}{Embaralhamento}
  Antes de dividir, embaralhe os dados com uma \textit{seed} fixa.
  Garante que as classes estejam distribuídas de forma similar nos três conjuntos.
\end{block}
\end{columns}
\end{frame}

% ── PASSO 5: MODELO ─────────────────────────────────────────
\begin{frame}{Passo 5 — Treinar o modelo}
\begin{columns}
\column{0.5\textwidth}
\textbf{O modelo: MLP (rede neural densa)}
\begin{itemize}
  \item Camadas ocultas com ativação \textbf{ReLU}
  \item Camada de saída com \textbf{Sigmoid} $\to$ probabilidade
  \item Cada camada aprende uma representação mais abstrata
\end{itemize}
\vspace{6pt}
\textbf{Função de perda: BCE}
\[
  \mathcal{L} = -\frac{1}{n}\sum_{i=1}^{n}
  \bigl[ y_i \log\hat{p}_i + (1-y_i)\log(1-\hat{p}_i) \bigr]
\]
\column{0.45\textwidth}
\textbf{Arquitetura usada:}\\[6pt]
\begin{center}
\begin{tabular}{ll}
  \toprule
  Camada & Operação \\
  \midrule
  Entrada  & 2 features \\
  Oculta 1 & Dense(16) + ReLU \\
  Oculta 2 & Dense(16) + ReLU \\
  Saída    & Dense(1)  + Sigmoid \\
  \bottomrule
\end{tabular}
\end{center}
\end{columns}
\end{frame}

% ── OTIMIZAÇÃO ──────────────────────────────────────────────
\begin{frame}{Como o modelo aprende}
\begin{columns}
\column{0.5\textwidth}
\textbf{Gradiente descendente estocástico (SGD):}
\begin{enumerate}
  \item Passa um mini-batch pelo modelo (\textit{forward})
  \item Calcula o gradiente da perda (\textit{backward})
  \item Atualiza os pesos na direção oposta ao gradiente:
  \[
    W \leftarrow W - \eta \cdot \nabla_W \mathcal{L}
  \]
  \item Repete por épocas
\end{enumerate}
\column{0.45\textwidth}
\begin{block}{O que é uma época?}
  Uma passagem completa pelo conjunto de treino.
  Com mini-batches, a cada época o modelo vê todos os exemplos em pequenos grupos.
\end{block}
\vspace{6pt}
\begin{block}{Por que mini-batches?}
  Equilíbrio entre estabilidade (batch completo)
  e velocidade/ruído benéfico (amostras individuais).
\end{block}
\end{columns}
\end{frame}

% ── VIÉS, VARIÂNCIA E GENERALIZAÇÃO ─────────────────────────
\begin{frame}{Viés, Variância e Generalização}
\begin{center}
\begin{tabular}{lcc}
  \toprule
  & \textbf{Treino} & \textbf{Validação / Teste} \\
  \midrule
  \textbf{Underfitting} (viés alto)   & Erro alto  & Erro alto \\
  \textbf{Situação ideal}              & Erro baixo & Erro baixo \\
  \textbf{Overfitting} (variância alta)& Erro baixo & Erro alto \\
  \bottomrule
\end{tabular}
\end{center}
\vfill
\begin{columns}
\column{0.32\textwidth}
\begin{alertblock}{Underfitting}
  Modelo muito simples. Não captura os padrões nos dados.\\
  \textit{Solução: mais capacidade, mais épocas, menos regularização.}
\end{alertblock}
\column{0.32\textwidth}
\begin{exampleblock}{Generalização}
  O modelo aprende os padrões do problema — e funciona em dados novos.\\
  \textit{Objetivo real de todo pipeline.}
\end{exampleblock}
\column{0.32\textwidth}
\begin{alertblock}{Overfitting}
  Modelo memoriza o treino. Falha em dados novos.\\
  \textit{Solução: mais dados, regularização, dropout, early stopping.}
\end{alertblock}
\end{columns}
\end{frame}

% ── DIAGNÓSTICO ─────────────────────────────────────────────
\begin{frame}{Diagnosticar pelo histórico de treino}
\begin{columns}
\column{0.5\textwidth}
\textbf{Leitura das curvas de loss:}
\begin{itemize}
  \item \textcolor{blue}{loss\_treino} cai: o modelo está aprendendo
  \item \textcolor{orange}{loss\_valid} cai junto: generalização boa
  \item \textcolor{orange}{loss\_valid} para e começa a subir: \textbf{overfitting}
  \item Ambas altas e estagnadas: \textbf{underfitting}
\end{itemize}
\vspace{8pt}
\textbf{Gap treino–validação:}
\[
  \Delta = \mathcal{L}_{\text{valid}} - \mathcal{L}_{\text{treino}}
\]
$\Delta$ pequeno $\to$ boa generalização \quad $\Delta$ grande $\to$ overfitting
\column{0.45\textwidth}
\begin{block}{Early stopping}
  Interrompe o treino quando a loss de validação não melhora por $p$ épocas consecutivas (\textit{paciência}).\\[4pt]
  \begin{itemize}
    \item Evita overfitting sem acertar o número de épocas na mão
    \item O melhor checkpoint é o de menor loss de validação
  \end{itemize}
\end{block}
\end{columns}
\end{frame}

% ── PASSO 6: AVALIAÇÃO ──────────────────────────────────────
\begin{frame}{Passo 6 — Avaliar e analisar erros}
\begin{columns}
\column{0.5\textwidth}
\textbf{Avaliar sempre no conjunto de teste — uma única vez.}\\[6pt]
\textbf{Relatório completo:}
\begin{itemize}
  \item Loss (BCE) — qualidade probabilística
  \item Acurácia — proporção de acertos
  \item Precisão — confiança ao prever positivo
  \item Recall — cobertura dos positivos reais
  \item F1-Score — balanço precisão/recall
  \item Matriz de confusão — mapa dos erros
\end{itemize}
\column{0.45\textwidth}
\begin{block}{FP vs FN — o custo depende do domínio}
  \begin{itemize}
    \item \textbf{Spam}: FP = e-mail legítimo na lixeira (incomoda)
    \item \textbf{Fraude}: FN = transação fraudulenta ignorada (caro)
    \item \textbf{Diagnóstico}: FN = doença não detectada (grave)
  \end{itemize}
\end{block}
\vspace{4pt}
\begin{alertblock}{Ciclo de análise}
  Avaliar $\to$ entender os erros $\to$ revisar dados ou modelo $\to$ retreinar
\end{alertblock}
\end{columns}
\end{frame}

% ── PASSO 7: PRODUÇÃO ───────────────────────────────────────
\begin{frame}{Passo 7 — Empacotar e colocar em produção}
\begin{columns}
\column{0.5\textwidth}
\textbf{O que vai para produção:}
\begin{itemize}
  \item Os \textbf{pesos e vieses} do modelo treinado
  \item Os \textbf{parâmetros do pré-processamento} ($\mu$, $\sigma$)
  \item \textbf{Não} os dados de treino
\end{itemize}
\vspace{6pt}
\textbf{Por que empacotar junto?}
\begin{itemize}
  \item Um dado novo em produção passa \textbf{exatamente} pelo mesmo pré-processamento do treino
  \item Separar os dois é garantia de inconsistência futura
\end{itemize}
\column{0.45\textwidth}
\begin{block}{Inferência em produção}
  \begin{enumerate}
    \item Recebe features brutas $X$
    \item Aplica pré-processamento: $X' = (X - \mu)/\sigma$
    \item Passa pelo modelo: $\hat{p} = \text{MLP}(X')$
    \item Decide: $\hat{y} = \mathbb{1}[\hat{p} \geq 0{,}5]$
    \item Retorna $\hat{p}$ e $\hat{y}$
  \end{enumerate}
\end{block}
\end{columns}
\end{frame}

% ── PRÓXIMOS PASSOS ─────────────────────────────────────────
\begin{frame}{Após o primeiro modelo: o que vem depois}
\begin{columns}
\column{0.48\textwidth}
\textbf{Melhorar o modelo:}
\begin{itemize}
  \item Busca de hiperparâmetros (learning rate, arquitetura)
  \item Regularização (L2, Dropout)
  \item Mais dados ou técnicas de aumento de dados
  \item Calibração de probabilidade
\end{itemize}
\column{0.48\textwidth}
\textbf{Sustentabilidade em produção:}
\begin{itemize}
  \item Monitoramento de \textit{data drift}
  \item Versionamento de modelos
  \item Re-treino periódico
  \item Testes automatizados do pipeline
\end{itemize}
\end{columns}
\vfill
\begin{block}{O pipeline é um ciclo, não uma linha reta}
  Dados $\to$ Modelo $\to$ Produção $\to$ Feedback $\to$ Novos Dados $\to$ \ldots
\end{block}
\end{frame}

% ── REFERÊNCIAS ─────────────────────────────────────────────
\begin{frame}{Referências}
\footnotesize
\begin{itemize}
  \item GOODFELLOW, I.; BENGIO, Y.; COURVILLE, A.
        \textit{Deep Learning}. MIT Press, 2016.
        Disponível em: \url{http://www.deeplearningbook.org}

  \item RODRIGUES, M. G. A.
        \textit{Redes Neurais Artificiais: Teoria e Aplicações}.
        TCC — Universidade de Brasília, 2023.
        \url{https://bdm.unb.br/bitstream/10483/36835/1/2023\_MatheusGabrielAlvesRodrigues\_tcc.pdf}

  \item IBM.
        \textit{What is a Machine Learning Pipeline?}
        IBM Think, 2024.
        \url{https://www.ibm.com/think/topics/machine-learning-pipeline}

  \item ALURA.
        \textit{Pipeline de IA: o que é e como funciona}.
        Alura, 2024.
        \url{https://www.alura.com.br/artigos/pipeline-de-ia}
\end{itemize}
\vfill
\begin{center}
  \small Código-fonte completo: \url{https://github.com/alexand7e/Pipeline-IA-do-Zero}
\end{center}
\end{frame}

% ── ENCERRAMENTO ────────────────────────────────────────────
\begin{frame}
\begin{center}
  \vfill
  {\Large \textbf{Obrigado!}}\\[12pt]
  \textit{Perguntas?}\\[20pt]
  \hrule
  \vspace{10pt}
  \small
  Alexandre Barros\\
  \href{https://alexand7e.dev.br}{alexand7e.dev.br}\\[4pt]
  \href{https://github.com/alexand7e/Pipeline-IA-do-Zero}{\texttt{github.com/alexand7e/Pipeline-IA-do-Zero}}
  \vfill
\end{center}
\end{frame}

\end{document}
